\documentclass[11pt,a4paper]{article}
\usepackage[utf8]{inputenc}
\usepackage[T1]{fontenc}
\usepackage[french]{babel}
\usepackage[a4paper,top=2cm,bottom=2.1cm,left=2.05cm,right=2.05cm]{geometry}
\usepackage{amsmath,amssymb}
\usepackage{array}
\usepackage{booktabs}
\usepackage{enumitem}
\usepackage{eurosym}
\usepackage{fancyhdr}
\usepackage{hyperref}
\usepackage{inconsolata}
\usepackage{listings}
\usepackage{microtype}
\usepackage{tabularx}
\usepackage[most]{tcolorbox}
\usepackage{titlesec}
\usepackage{xcolor}

\definecolor{ink}{HTML}{202124}
\definecolor{muted}{HTML}{5F6368}
\definecolor{line}{HTML}{DADCE0}
\definecolor{paper}{HTML}{F8F9FA}
\definecolor{accent}{HTML}{8A1538}
\definecolor{accentsoft}{HTML}{F8E8EE}
\definecolor{green}{HTML}{0B6B57}
\definecolor{greensoft}{HTML}{E6F4EF}
\definecolor{amber}{HTML}{9A5B00}
\definecolor{ambersoft}{HTML}{FFF4D6}
\definecolor{blue}{HTML}{245B8E}
\definecolor{bluesoft}{HTML}{EAF2FA}
\definecolor{violet}{HTML}{5C4D7D}
\definecolor{violetsoft}{HTML}{F1EEF8}

\hypersetup{
  colorlinks=true,
  linkcolor=accent,
  urlcolor=green
}

\setlength{\parindent}{0pt}
\setlength{\parskip}{5pt}
\setlist[itemize]{leftmargin=1.35em,itemsep=2pt,topsep=2pt}
\setlist[enumerate]{leftmargin=1.55em,itemsep=3pt,topsep=3pt}
\setlist[enumerate,1]{label=\textbf{\alph*.}}

\pagestyle{fancy}
\fancyhf{}
\fancyhead[L]{Python pour économistes}
\fancyhead[R]{Université de Lille -- L2 Économie}
\fancyfoot[C]{\thepage}
\renewcommand{\headrulewidth}{0.35pt}
\setlength{\headheight}{14.5pt}

\titleformat{\section}
  {\Large\bfseries\color{accent}}
  {}{0pt}{}
  [\vspace{-0.25em}\color{line}\titlerule]
\titleformat{\subsection}
  {\large\bfseries\color{ink}}
  {}{0pt}{}

\lstdefinelanguage{terminal}{
  morekeywords={ls,cd,pwd,mkdir,python,python3,code,dir,type,echo},
  sensitive=true
}

\lstset{
  basicstyle=\ttfamily\small,
  backgroundcolor=\color{paper},
  frame=single,
  framerule=0.35pt,
  rulecolor=\color{line},
  language=Python,
  keywordstyle=\color{accent}\bfseries,
  commentstyle=\color{muted},
  stringstyle=\color{green},
  showstringspaces=false,
  breaklines=true,
  columns=fullflexible,
  keepspaces=true,
  upquote=true,
  literate=
   {à}{{\`a}}1 {â}{{\^a}}1 {ä}{{\"a}}1
   {ç}{{\c c}}1
   {é}{{\'e}}1 {è}{{\`e}}1 {ê}{{\^e}}1 {ë}{{\"e}}1
   {î}{{\^\i}}1 {ï}{{\"\i}}1
   {ô}{{\^o}}1 {ö}{{\"o}}1
   {ù}{{\`u}}1 {û}{{\^u}}1 {ü}{{\"u}}1
   {À}{{\`A}}1 {Â}{{\^A}}1 {Ä}{{\"A}}1
   {Ç}{{\c C}}1
   {É}{{\'E}}1 {È}{{\`E}}1 {Ê}{{\^E}}1 {Ë}{{\"E}}1
   {Î}{{\^I}}1 {Ï}{{\"I}}1
   {Ô}{{\^O}}1 {Ö}{{\"O}}1
   {Ù}{{\`U}}1 {Û}{{\^U}}1 {Ü}{{\"U}}1
   {œ}{{\oe}}1 {Œ}{{\OE}}1
   {–}{{--}}1 {—}{{---}}1
}

\newcommand{\code}[1]{\texttt{#1}}
\newcommand{\pp}{\,points de pourcentage}

\newcounter{exercise}
\newcounter{extension}
\newcounter{logic}

\newcommand{\workbooktitle}[4]{%
  \setcounter{exercise}{0}%
  \setcounter{extension}{0}%
  \setcounter{logic}{0}%
  \fancyhead[L]{#1 -- #2}%
  {\Large\bfseries\color{ink}#1 -- #2\par}
  \vspace{0.2em}
  {\color{muted}Python pour économistes -- Université de Lille, L2 Économie\par}
  \vspace{0.65em}
  {\color{accent}\rule{\linewidth}{0.8pt}}
  \vspace{0.5em}
  \begin{routebox}{Question fil rouge}
  #3
  \end{routebox}
  \begin{deliverablebox}{Livrables}
  #4
  \end{deliverablebox}
  \begin{methodbox}{Méthode de travail}
  \begin{tabularx}{\linewidth}{@{}>{\bfseries\color{blue}}p{0.18\linewidth}X@{}}
  Lire & identifier les données, la question et le livrable attendu.\\
  Décomposer & écrire les étapes en pseudo-code avant de coder.\\
  Exécuter & lancer le script complet depuis le bon dossier.\\
  Vérifier & contrôler les sorties, les fichiers créés et l’interprétation.
  \end{tabularx}
  \end{methodbox}
}

\newtcolorbox{routebox}[1]{
  enhanced,breakable,
  title={#1},
  colback=paper,
  colframe=line,
  coltitle=ink,
  fonttitle=\bfseries,
  boxrule=0.5pt,
  arc=2pt,
  left=8pt,right=8pt,top=7pt,bottom=7pt,
  borderline west={2.4pt}{0pt}{accent}
}

\newtcolorbox{deliverablebox}[1]{
  enhanced,breakable,
  title={#1},
  colback=greensoft,
  colframe=green,
  coltitle=ink,
  fonttitle=\bfseries,
  boxrule=0.6pt,
  arc=2pt,
  left=8pt,right=8pt,top=7pt,bottom=7pt,
  borderline west={2.4pt}{0pt}{accent}
}


\newtcolorbox{methodbox}[1]{
  enhanced,breakable,
  title={#1},
  colback=bluesoft,
  colframe=blue,
  coltitle=ink,
  fonttitle=\bfseries,
  boxrule=0.55pt,
  arc=2pt,
  left=8pt,right=8pt,top=6pt,bottom=6pt,
  before skip=7pt,after skip=9pt,
  borderline west={2.4pt}{0pt}{blue}
}

\newtcolorbox{pathwaybox}[1]{
  enhanced,breakable,
  title={#1},
  colback=violetsoft,
  colframe=violet,
  coltitle=ink,
  fonttitle=\bfseries,
  boxrule=0.55pt,
  arc=2pt,
  left=8pt,right=8pt,top=6pt,bottom=6pt,
  before skip=8pt,after skip=8pt,
  borderline west={2.4pt}{0pt}{violet}
}

\newenvironment{pseudocodeexercise}[1]{%
  \refstepcounter{logic}%
  \begin{tcolorbox}[
    enhanced,breakable,
    title={Logique \thelogic{} -- #1 \hfill\scriptsize PSEUDO-CODE},
    colback=bluesoft,
    colframe=blue,
    coltitle=ink,
    fonttitle=\bfseries,
    boxrule=0.6pt,
    arc=2pt,
    left=8pt,right=8pt,top=6pt,bottom=6pt,
    before skip=8pt,after skip=8pt,
    borderline west={2.6pt}{0pt}{blue}
  ]%
}{\end{tcolorbox}}

\newenvironment{mandatoryexercise}[1]{%
  \refstepcounter{exercise}%
  \begin{tcolorbox}[
    enhanced,breakable,
    title={Exercice \theexercise{} -- #1 \hfill\scriptsize OBLIGATOIRE},
    colback=white,
    colframe=accent,
    coltitle=ink,
    fonttitle=\bfseries,
    boxrule=0.7pt,
    arc=2pt,
    left=8pt,right=8pt,top=7pt,bottom=7pt,
    before skip=9pt,after skip=9pt,
    borderline west={2.4pt}{0pt}{accent}
  ]%
}{\end{tcolorbox}}

\newenvironment{extensionexercise}[1]{%
  \refstepcounter{extension}%
  \begin{tcolorbox}[
    enhanced,breakable,
    title={Extension \theextension{} -- #1},
    colback=ambersoft,
    colframe=amber,
    coltitle=ink,
    fonttitle=\bfseries,
    boxrule=0.55pt,
    arc=2pt,
    left=8pt,right=8pt,top=7pt,bottom=7pt,
    before skip=8pt,after skip=8pt,
    borderline west={2.4pt}{0pt}{amber}
  ]%
}{\end{tcolorbox}}

\newtcolorbox{checkpointbox}[1]{
  enhanced,breakable,
  title={#1},
  colback=accentsoft,
  colframe=accent,
  coltitle=ink,
  fonttitle=\bfseries,
  boxrule=0.55pt,
  arc=2pt,
  left=8pt,right=8pt,top=7pt,bottom=7pt,
  borderline west={2.4pt}{0pt}{accent}
}


\begin{document}

\workbooktitle
  {TP5}
  {Mini-projet pandas avec Our World in Data}
  {Comment importer un CSV OWID, produire des statistiques descriptives simples, tracer un graphique et rédiger une interprétation économique ?}
  {Un fichier \code{tp5\_nom\_prenom.py}, une figure \code{out/gdp\_top10.png}, un fichier \code{out/gdp\_selection.csv} et une interprétation de six à huit lignes.}

\begin{checkpointbox}{Parcours obligatoire}
Les exercices 1 à 9 sont obligatoires. Ce TP est le mini-projet final : il combine chemins de fichiers, \code{pandas}, statistiques descriptives simples, graphique et interprétation.
\end{checkpointbox}

\begin{pathwaybox}{Pathways : travailler avec un CSV OWID}
Le TP lit directement un CSV publié par Our World in Data. Si la connexion internet ne fonctionne pas en salle, télécharger le CSV avec un navigateur et placer le fichier dans un dossier \code{data}.
\begin{lstlisting}[language=terminal]
https://ourworldindata.org/grapher/gdp-per-capita-worldbank.csv
\end{lstlisting}
\end{pathwaybox}

\begin{pseudocodeexercise}{Pipeline pandas final}
\begin{lstlisting}
# Objectif : mini-projet descriptif avec un CSV OWID
# 1. Lire le CSV
# 2. Inspecter lignes, colonnes et types
# 3. Renommer les colonnes utiles
# 4. Filtrer une annee et quelques pays
# 5. Calculer moyenne, mediane, min, max
# 6. Tracer un graphique lisible
# 7. Exporter un tableau propre
# 8. Rediger une interpretation prudente
\end{lstlisting}
\end{pseudocodeexercise}

\section*{A. Importer et inspecter}

\begin{mandatoryexercise}{Préparer le script}
Créer l'arborescence suivante :
\begin{lstlisting}[language=terminal]
TP5/
  tp5_nom_prenom.py
  out/
\end{lstlisting}
Dans le script, créer le dossier \code{out}.
\begin{lstlisting}
from pathlib import Path

out_dir = Path("out")
out_dir.mkdir(exist_ok=True)
\end{lstlisting}
\end{mandatoryexercise}

\begin{mandatoryexercise}{Lire le CSV OWID}
Lire le fichier depuis l'URL OWID.
\begin{lstlisting}
import pandas as pd

url = "https://ourworldindata.org/grapher/gdp-per-capita-worldbank.csv"
df = pd.read_csv(url)

print(df.head())
print(df.shape)
print(df.columns)
\end{lstlisting}
Ajouter un commentaire indiquant le nombre de lignes et de colonnes.
\end{mandatoryexercise}

\begin{mandatoryexercise}{Renommer proprement}
Le nom de la colonne de valeur peut être long. On le récupère automatiquement.
\begin{lstlisting}
value_col = df.columns[-1]

df = df.rename(columns={
    "Entity": "country",
    "Code": "code",
    "Year": "year",
    value_col: "gdp_pc"
})
\end{lstlisting}
Afficher \code{df[["country", "year", "gdp\_pc"]].head()}.
\end{mandatoryexercise}

\section*{B. Filtrer et décrire}

\begin{mandatoryexercise}{Sélectionner une année récente}
Créer un tableau \code{latest} avec la dernière année disponible dans le fichier, puis retirer les valeurs manquantes.
\begin{lstlisting}
latest_year = df["year"].max()
latest = df.loc[df["year"] == latest_year].copy()
latest = latest.dropna(subset=["gdp_pc"])
print(latest_year)
print(latest.shape)
\end{lstlisting}
\end{mandatoryexercise}

\begin{mandatoryexercise}{Sélectionner quelques pays}
Créer une liste de pays et filtrer le tableau.
\begin{lstlisting}
pays_selection = [
    "France", "Germany", "Spain", "Poland",
    "United States", "China", "India", "Brazil"
]
selection = latest.loc[latest["country"].isin(pays_selection)].copy()
print(selection[["country", "gdp_pc"]])
\end{lstlisting}
\end{mandatoryexercise}

\begin{mandatoryexercise}{Statistiques descriptives simples}
Sur \code{selection}, calculer :
\begin{enumerate}
  \item moyenne ;
  \item médiane ;
  \item minimum ;
  \item maximum ;
  \item proportion de pays au-dessus de la moyenne de la sélection.
\end{enumerate}
On peut utiliser \code{mean()}, \code{median()}, \code{min()} et \code{max()}.
\end{mandatoryexercise}

\section*{C. Graphique et export}

\begin{mandatoryexercise}{Graphique du top 10}
Créer un graphique en barres des dix pays les plus élevés dans \code{latest}.
\begin{lstlisting}
import matplotlib.pyplot as plt

top10 = latest.sort_values("gdp_pc", ascending=False).head(10)
ax = top10.plot(kind="bar", x="country", y="gdp_pc", legend=False)
ax.set_xlabel("Pays")
ax.set_ylabel("PIB par habitant")
ax.set_title(f"Top 10 du PIB par habitant - {latest_year}")
plt.xticks(rotation=45, ha="right")
plt.tight_layout()
plt.savefig(out_dir / "gdp_top10.png", dpi=300)
\end{lstlisting}
\end{mandatoryexercise}

\begin{mandatoryexercise}{Exporter un tableau propre}
Exporter \code{selection} dans \code{out/gdp\_selection.csv} avec seulement \code{country}, \code{code}, \code{year}, \code{gdp\_pc}.
\begin{lstlisting}
selection[["country", "code", "year", "gdp_pc"]].to_csv(
    out_dir / "gdp_selection.csv",
    index=False
)
\end{lstlisting}
\end{mandatoryexercise}

\begin{mandatoryexercise}{Interprétation finale}
Ajouter six à huit lignes de commentaire :
\begin{enumerate}
  \item année étudiée ;
  \item pays le plus élevé dans la sélection ;
  \item pays le plus faible dans la sélection ;
  \item moyenne et médiane de la sélection ;
  \item comparaison de la France à la moyenne ;
  \item limite : choix des pays ;
  \item limite : le PIB par habitant ne mesure pas toutes les dimensions du développement.
\end{enumerate}
\end{mandatoryexercise}

\section*{D. Extension facultative}

\begin{extensionexercise}{Pour les étudiants plus rapides}
Choisir au moins deux questions si le mini-projet est terminé.
\begin{enumerate}
  \item Tracer uniquement les pays de la sélection.
  \item Comparer la France à l'Allemagne, aux États-Unis et à la Chine.
  \item Créer une colonne \code{above\_selection\_mean}.
  \item Exporter le \code{top10} dans \code{out/gdp\_top10.csv}.
  \item Tracer l'évolution du PIB par habitant de la France depuis 1990.
  \item Ajouter le monde \code{"World"} si présent dans le fichier.
  \item Remplacer \code{latest\_year} par une année choisie manuellement.
  \item Calculer le rapport entre le maximum et le minimum de la sélection.
  \item Importer aussi l'espérance de vie depuis \url{https://ourworldindata.org/grapher/life-expectancy.csv}.
  \item Extension avancée : fusionner PIB par habitant et espérance de vie pour une même année.
\end{enumerate}
\end{extensionexercise}

\end{document}
